\section{一次完整的前向与反向传播}

反向传播不是参数更新算法。它只做一件事：给定一个标量损失，利用计算图中的局部导数，高效算出损失对所有参数的梯度。梯度下降、Adam 或其他优化器随后才决定怎样使用这些梯度。

\subsection{先用一个标量网络看清链式法则}

考虑

\[
z=w_1x+b_1,\qquad
h=\operatorname{ReLU}(z),\qquad
o=w_2h+b_2,
\]

\[
L=\frac12(o-y)^2.
\]

取 \(x=2,w_1=0.5,b_1=0,w_2=-1,b_2=0,y=1\)。前向得到

\[
z=1,\quad h=1,\quad o=-1,\quad L=2.
\]

从损失开始向后：

\[
\frac{\partial L}{\partial o}=o-y=-2,
\]

\[
\frac{\partial L}{\partial w_2}
=\frac{\partial L}{\partial o}
\frac{\partial o}{\partial w_2}
=-2h=-2,
\qquad
\frac{\partial L}{\partial b_2}=-2.
\]

传回隐藏单元：

\[
\frac{\partial L}{\partial h}
=\frac{\partial L}{\partial o}w_2=2.
\]

由于 \(z>0\)，ReLU 导数为 1，所以

\[
\frac{\partial L}{\partial z}=2,
\quad
\frac{\partial L}{\partial w_1}=2x=4,
\quad
\frac{\partial L}{\partial b_1}=2.
\]

前向传播保存 \(x,z,h,o\) 等局部计算所需信息；反向传播从标量 1 或损失导数出发，把上游敏感度乘以局部导数。若一个变量沿多条路径影响损失，各路径贡献必须相加，而不是覆盖。

\subsection{Batch 矩阵形式只是同一逻辑的压缩}

采用``样本按行''的约定。设 batch 大小为 \(B\)，各层宽度为 \(d_0,d_1,d_2\)，则 \(X\in\R^{B\times d_0}\)，\(W_1\in\R^{d_0\times d_1}\)，\(W_2\in\R^{d_1\times d_2}\)，\(Z_1,H\in\R^{B\times d_1}\)，\(Z_2\in\R^{B\times d_2}\)。前向计算为

\[
Z_1=XW_1+\mathbf1b_1^\top,
\qquad H=\phi(Z_1),
\]

\[
Z_2=HW_2+\mathbf1b_2^\top,
\qquad
L=\frac1B\sum_{i=1}^B\ell(Z_{2i},y_i).
\]

若

\[
G_2=\frac{\partial L}{\partial Z_2},
\]

逐样本看，\(Z_2\) 的第 \(i\) 行只参与第 \(i\) 个样本的损失，该样本对 \(\nabla_{W_2}L\) 的贡献是外积 \(h_ig_{2i}^{\top}\)（\(h_i^{\top}\)、\(g_{2i}^{\top}\) 分别为 \(H\) 与 \(G_2\) 的第 \(i\) 行）。对 \(i\) 求和恰是矩阵乘积，于是第二层梯度为

\[
\nabla_{W_2}L=H^\top G_2,
\qquad
\nabla_{b_2}L=G_2^\top\mathbf1.
\]

传回隐藏层（\(\odot\) 表示逐元素乘积）：

\[
G_1
=(G_2W_2^\top)\odot\phi'(Z_1),
\]

于是

\[
\nabla_{W_1}L=X^\top G_1,
\qquad
\nabla_{b_1}L=G_1^\top\mathbf1.
\]

每个参数梯度 shape 必须与参数相同。bias 在前向被广播到 batch 中每一行，反向就要沿 batch 轴求和。若损失取 batch mean，\(1/B\) 应已经出现在 \(G_2\) 中；漏掉或重复这个因子会让学习率随 batch size 意外改变。

对 softmax cross-entropy，

\[
G_2=\frac{P-Y}{B},
\]

其中 \(P\) 是对 \(Z_2\) 逐行做 softmax 得到的概率矩阵，\(Y\) 是对应的 one-hot 标签矩阵；这正是上一节得到的 \(p-y\) 信号的 batch 形式。代入上式即可得到两层分类器的完整参数梯度。真实框架的 reverse-mode automatic differentiation 正是在更大计算图上执行同类 vector--Jacobian product，而不是显式构造每个巨大 Jacobian。

\subsection{梯度检查是局部证据}

对方向 \(v\)，中心有限差分

\[
\frac{
L(\theta+\varepsilon v)
-L(\theta-\varepsilon v)
}{2\varepsilon}
\approx
\nabla L(\theta)^\top v
\]

其中 Hessian 项在中心差分中相消，截断误差为 \(O(\varepsilon^2)\)。取若干随机方向即可核对整组梯度。使用 double precision、小网络和不过大也不过小的 \(\varepsilon\)；避开 ReLU 恰好在零点的不可微位置。检查相对误差，并单独测试 batch size 1、bias broadcast、mask 和共享参数。

梯度检查通过只说明当前点附近的导数实现大致正确，不说明模型、损失、数据切分或优化方案正确。in-place 修改缓存张量、无意 detach、transpose 约定混乱和多路径梯度未累加，都是自动微分图中常见错误。

\subsubsection{手算后再交给框架}

把标量例子改为两个隐藏单元，并让两个单元的输出都进入同一个最终节点。先手算每条路径对 \(x\) 的贡献并相加，再用自动微分核对。随后实现上面的 batch 网络，对一个随机方向做有限差分，而不是逐元素检查所有参数。解释为什么 directional check 能有效发现 transpose 与 reduction 错误。

矩阵微积分与计算图的通用形式见``计算图上的反向传播与自动微分''。
